\section{Разбиение области IAPWS-IF97 на регионы}

Рабочая область IF97 разделена на регионы,
в каждом из которых применяются собственные уравнения
и собственные нормировки безразмерных переменных.
Такое разбиение отражает физические различия
между сжатой жидкостью, перегретым паром и околокритической областью.
Одна универсальная аналитическая формула для всей области
в промышленной постановке не обеспечивает одновременно
точность, устойчивость и вычислительную эффективность.

Следовательно, для любого маршрута расчёта
определение региона является обязательным этапом.
Только после классификации состояния корректно
выбирается модель и вычисляются итоговые свойства.

\textbf{Краткая характеристика стандарта IAPWS-IF97.}
IF97 представляет собой промышленную региональную формулировку
для чистой воды и водяного пара,
ориентированную на инженерные расчёты,
где одновременно важны точность,
вычислительная скорость
и устойчивость на границах областей~\cite{iapws_if97_2012}.
Стандарт не задаёт одно универсальное уравнение
	для всей рабочей области,
	а использует набор специализированных аппроксимаций:
	основные уравнения для регионов~1--5,
	аналитическую зависимость для линии насыщения региона~4,
	а также часть обратных зависимостей для входных пар
	$(p, h)$ и $(p, s)$.
При этом IF97 имеет и ограничения:
он относится к чистой воде,
	не предназначен для солевых растворов и смесей,
	не покрывает все возможные обратные постановки
	в аналитической форме,
	а для региона~5 не учитывает эффекты термической диссоциации.
Именно поэтому в вычислительном ядре
аналитические соотношения IF97
дополняются явной валидацией,
маршрутизацией по регионам
и численными процедурами на сложных участках области.

Геометрия рабочей области стандарта наглядно показана на рисунке~\ref{fig:if97_regions}.
Из него видно, что регионы разделены не произвольно, а по физически значимым границам:
линиям насыщения, критической области и высоким температурам,
поэтому выбор формулы в ядре всегда связан с положением точки на этой диаграмме.

\begin{figure}[htbp]
    \centering
    \begin{tikzpicture}
        
        % ЛЕВЫЙ ГРАФИК: Регионы 1, 2, 3 и 4
        \begin{axis}[
            name=ax1,
            width=10cm, height=8cm, % Немного сузили, чтобы точно влезло в поля
            xmin=250, xmax=1150,
            ymin=0, ymax=110,
            axis y line=left,
            axis x line=bottom,
            ylabel={Давление $p$, МПа},
            xlabel={Температура $T$, К},
            xtick={273.15, 623.15, 1073.15},
            ytick={0, 50, 100}, % Оставили только главные метки
            extra y ticks={22.06}, % Добавили критическое давление отдельно
            extra y tick labels={22{,}06},
            grid=both,
            grid style={dashed, gray!30},
            enlargelimits=false,
            ticklabel style={/pgf/number format/use comma},
            clip=false % Отключаем обрезку элементов, выходящих за рамки
        ]
        
        % Регион 4: Линия насыщения (по реперным точкам IF97)
        \addplot [orange, ultra thick, smooth] coordinates {
            (273.15, 0.000611)
            (373.15, 0.1013)
            (473.15, 1.554)
            (573.15, 8.581)
            (647.096, 22.064)
        };
        
        % Граница между Регионами 2 и 3 (Уравнение B23)
        \addplot [domain=623.15:863.15, samples=50, thick, dashed] 
            {13.918837 + 0.001019297 * (x - 572.5446)^2};
            
        % Вертикальная граница между Регионами 1 и 3
        \draw[thick, dashed] (axis cs:623.15, 16.53) -- (axis cs:623.15, 100);
        
        % Вертикальная граница между Регионами 2 и 5
        \draw[thick] (axis cs:1073.15, 0) -- (axis cs:1073.15, 100);
        
        % Верхняя граница графика (p = 100 МПа)
        \draw[thick] (axis cs:273.15, 100) -- (axis cs:1073.15, 100);
        
        % Подписи регионов
        \node at (axis cs: 440, 60) [font=\large] {Регион 1};
        \node at (axis cs: 870, 30) [font=\large] {Регион 2};
        \node at (axis cs: 720, 80) [font=\large] {Регион 3};
        
        % Выноска для Региона 4 (сделали аккуратнее)
        \draw[<-, thick] (axis cs: 573.15, 8.5) -- (axis cs: 750, 15) 
            node[anchor=west, font=\large] {Регион 4};
            
        \end{axis}
        
        % ПРАВЫЙ ГРАФИК: Регион 5 (Разрыв оси)
        \begin{axis}[
            name=ax2,
            at={(ax1.south east)},
            anchor=south west, % Жестко выравниваем левый нижний угол
            xshift=0.8cm, % Увеличили отступ для значка разрыва
            width=3.5cm, height=8cm,
            xmin=2173.15, xmax=2373.15,
            ymin=0, ymax=110,
            axis y line=none, 
            axis x line=bottom,
            xtick={2273.15},
            grid=both,
            grid style={dashed, gray!30},
            enlargelimits=false,
            ticklabel style={/pgf/number format/use comma},
            clip=false % Важно! Чтобы надписи не обрезались правым краем
        ]
        
        % Верхняя и правая границы Региона 5
        \draw[thick] (axis cs:2173.15, 50) -- (axis cs:2273.15, 50);
        \draw[thick] (axis cs:2273.15, 0) -- (axis cs:2273.15, 50);
        
        \node at (axis cs: 2223, 25) [font=\large] {Регион 5};
        \node at (axis cs: 2223, 55) {50 МПа};
        
        \end{axis}
        
        % Отрисовка значка "разрыва" оси X между графиками
        \path (ax1.south east) -- (ax2.south west) coordinate[midway] (mid);
        \draw[thick] (mid) ++(-0.15cm, -0.2cm) -- ++(0.1cm, 0.4cm);
        \draw[thick] (mid) ++(0.05cm, -0.2cm) -- ++(0.1cm, 0.4cm);
        
    \end{tikzpicture}
    \caption{Диаграмма $p$--$T$ с границами регионов IAPWS-IF97}
    \label{fig:if97_regions_fixed}
\end{figure}

